\documentclass[12pt,a4paper]{article}
\usepackage[utf8]{inputenc}
\usepackage[ngerman]{babel}
\usepackage{amsmath}
\usepackage{geometry}
\usepackage{booktabs}
\usepackage{listings}
\usepackage{xcolor}
\usepackage{hyperref}

\geometry{margin=2.5cm}

\title{Skill-Extraktion in der Datenbank\\[0.5em]
\large Dokumentation der Skripte\\
\texttt{extract\_skills\_from\_cv\_pdf.py}, \\
\texttt{extract\_skills\_from\_long\_note.py}, \\
\texttt{job\_skills\_extraktion.py}}
\author{KI-unterstütztes BMS}
\date{\today}

\begin{document}

\maketitle
\tableofcontents
\newpage

\section{Übersicht}

Die drei Skripte
\begin{itemize}
    \item \texttt{src/PostreSQL/import/extract\_skills\_from\_cv\_pdf.py},
    \item \texttt{src/PostreSQL/import/extract\_skills\_from\_long\_note.py} und
    \item \texttt{src/PostreSQL/import/job\_skills\_extraktion.py}
\end{itemize}
verwenden eine gemeinsame Logik, um fachliche \emph{Skills} aus Freitexten zu erkennen und in der PostgreSQL-Datenbank abzulegen.

Zentrale Ziele:
\begin{itemize}
    \item Skills aus verschiedenen Quellen (PDF-Lebenslauf, Freitextfeld, Stellenbeschreibung) einheitlich erkennen.
    \item Eine zentrale Skill-Liste (\texttt{data/skills.csv}) als Referenz verwenden.
    \item Gefundene Skills strukturiert und normalisiert in den Tabellen \texttt{candidates} bzw. \texttt{jobs} speichern.
\end{itemize}

\section{Datenbasis: \texttt{data/skills.csv}}

Alle drei Skripte laden dieselbe CSV-Datei \texttt{data/skills.csv}. Erwartete Spalten:
\begin{itemize}
    \item \textbf{Bezeichnung}: kanonischer Name des Skills (z.B. ``Python'', ``Radiologie'').
    \item \textbf{Gängige Abkürzungen}: optionale, komma-separierte Liste von Synonymen und Abkürzungen
          (z.B. ``Py, Python3'').
\end{itemize}

Aus jeder Zeile der CSV wird eine Menge von Suchbegriffen aufgebaut:
\begin{enumerate}
    \item Start mit der Bezeichnung selbst.
    \item Falls vorhanden, Aufspaltung der Abkürzungen an Kommas und Trimmen der Einträge.
    \item Alle Begriffe werden in Kleinbuchstaben umgewandelt und Duplikate entfernt.
\end{enumerate}

Das Ergebnis ist eine Mapping-Struktur
\[
\texttt{skills\_map: \{canonical\_name $\to$ [suchbegriffe]\}},
\]
die die Grundlage für alle weiteren Suchen bildet.

\section{Matching-Logik: Exakte und unscharfe Suche}

Die Kernfunktion zur Erkennung von Skills in Freitexten heißt in allen Skripten sinngemäß
\texttt{extract\_skills\_from\_text}. Sie verwendet zwei Schritte:

\subsection{Exakte Worttreffer}

Zunächst wird versucht, jeden Suchbegriff als ganzes Wort im Text zu finden. Dazu wird der
Text in Kleinbuchstaben umgewandelt und mit einem regulären Ausdruck durchsucht:
\begin{itemize}
    \item Muster: \verb|\b<suchbegriff>\b| (Wortgrenzen vor und nach dem Begriff).
    \item Der Suchbegriff wird mit \verb|re.escape| maskiert, um Sonderzeichen korrekt zu behandeln.
\end{itemize}

Wird ein solcher exakter Treffer gefunden, wird der zugehörige kanonische Skill-Name zur
Ergebnismenge hinzugefügt.

\subsection{Unscharfe Suche (Fuzzy Matching)}

Falls kein exakter Treffer gefunden wird, kommt eine unscharfe Suche mit der Bibliothek
\texttt{fuzzywuzzy} zum Einsatz:
\begin{itemize}
    \item Für den gesamten Text (in Kleinbuchstaben) und den Suchbegriff wird
          \verb|fuzz.partial_ratio(text, term)| berechnet.
    \item Liegt der Score oberhalb eines Schwellwertes (Standard: \(85\)),
          gilt der Skill als gefunden.
\end{itemize}

So werden auch Skills erkannt, wenn z.B. Beugungen, Tippfehler oder leicht andere Schreibweisen
verwendet werden.

\section{Skript: \texttt{extract\_skills\_from\_long\_note.py}}

Dieses Skript extrahiert Skills aus der Spalte \texttt{long\_note} eines einzelnen Kandidaten.

\subsection{Ablauf}
\begin{enumerate}
    \item Skills aus \texttt{data/skills.csv} laden (siehe Abschnitt~\ref{sec:datenbasis}).
    \item Verbindung zur Datenbank aufbauen (\texttt{DB\_CONFIG}).
    \item Kandidaten-ID vom Nutzer abfragen.
    \item Den entsprechenden Datensatz aus der Tabelle \texttt{candidates} laden:
    \begin{itemize}
        \item Spalten: \texttt{id}, \texttt{first\_name}, \texttt{last\_name}, \texttt{long\_note}, \texttt{skills}.
    \end{itemize}
    \item Vorschau auf den Inhalt von \texttt{long\_note} anzeigen (max. 300 Zeichen).
    \item Mit \texttt{extract\_skills\_from\_text} die Skills aus \texttt{long\_note} extrahieren.
    \item Gefundene Skills alphabetisch sortieren und als kommaseparierte Liste zusammenfassen.
    \item Gefundene Skills anzeigen und den Nutzer um Bestätigung bitten.
    \item Bei Bestätigung: \texttt{UPDATE candidates SET skills = ... WHERE id = ...}.
\end{enumerate}

\section{Skript: \texttt{extract\_skills\_from\_cv\_pdf.py}}

Dieses Skript extrahiert Skills aus einem in der Spalte \texttt{cv\_pdf} (BYTEA) gespeicherten
PDF-Lebenslauf eines Kandidaten und schreibt das Ergebnis in \texttt{candidates.skills}.

\subsection{Zusätzliche Schritte}
\begin{itemize}
    \item Das PDF wird aus dem BYTEA-Feld \texttt{cv\_pdf} geladen.
    \item Mittels \texttt{PyPDF2.PdfReader} wird Seiteninhalt in Text umgewandelt.
    \item Der extrahierte Volltext dient als Eingabe für
          \texttt{extract\_skills\_from\_text}.
\end{itemize}

\subsection{Ablauf (vereinfacht)}
\begin{enumerate}
    \item Skills laden, DB-Verbindung herstellen.
    \item Kandidaten-ID vom Nutzer abfragen.
    \item Datensatz aus \texttt{candidates} inkl. \texttt{cv\_pdf} laden.
    \item PDF in Text umwandeln und kurze Vorschau anzeigen.
    \item Skills aus dem CV-Text extrahieren.
    \item Gefundene Skills anzeigen und Bestätigung einholen.
    \item Bei Bestätigung: Skills-Feld in \texttt{candidates} aktualisieren.
\end{enumerate}

\section{Skript: \texttt{job\_skills\_extraktion.py}}

Dieses Skript überträgt die gleiche Skill-Logik auf die Tabelle \texttt{jobs}. Statt eines
interaktiven Einzelfalls werden hier viele Jobs in einem Durchlauf verarbeitet.

\subsection{Zielspalte}

Die erkannten Skills werden in der Spalte \texttt{sonstiges\_anforderungen} der Tabelle
\texttt{jobs} gespeichert. Grundlage der Erkennung sind die beiden Freitextspalten
\texttt{job\_description} und \texttt{long\_note}.

\subsection{Filterkriterium}

Aktualisiert werden nur Datensätze, bei denen
\begin{itemize}
    \item mindestens eine der Spalten \texttt{job\_description} oder \texttt{long\_note}
          nicht leer ist, und
    \item \texttt{sonstiges\_anforderungen} aktuell leer oder \texttt{NULL} ist.
\end{itemize}

\subsection{Ablauf}
\begin{enumerate}
    \item Skills aus \texttt{data/skills.csv} laden.
    \item Verbindung zur Datenbank herstellen.
    \item Geeignete Jobs mit obigem Filter aus \texttt{jobs} laden
          (inkl. \texttt{id}, \texttt{position}, \texttt{department},
           \texttt{job\_description}, \texttt{long\_note}, \texttt{sonstiges\_anforderungen}).
    \item Für jeden Job wird der Text aus \texttt{job\_description} und
          \texttt{long\_note} zu einem gemeinsamen Freitext zusammengefügt.
    \item Mit \texttt{extract\_skills\_from\_text} werden Skills aus diesem Text extrahiert.
    \item Für jeden Job mit gefundenen Skills wird ein Update-Vorschlag gebaut:
          neue Skills als sortierte, kommaseparierte Liste.
    \item Für die ersten Einträge wird eine Vorschau (Alter Wert / Neuer Wert)
          auf der Konsole ausgegeben.
    \item Gesamtanzahl der geplanten Updates wird angezeigt.
    \item Nach Bestätigung durch den Nutzer werden alle geplanten Updates per
          \texttt{UPDATE jobs SET sonstiges\_anforderungen = ... WHERE id = ...}
          in der Datenbank ausgeführt.
\end{enumerate}

\section{Gemeinsamkeiten und Unterschiede}

\subsection{Gemeinsame Bausteine}

\begin{itemize}
    \item Verwendung derselben Skill-Datenbasis (\texttt{data/skills.csv}).
    \item Gemeinsame Matching-Logik mit exakten Worttreffern und Fuzzy Matching.
    \item Speicherung der Ergebnisse als normalisierte, alphabetisch sortierte
          Liste von Skill-Namen im jeweiligen Ziel-Feld.
    \item Interaktive Bestätigung vor dem Schreiben in die Datenbank.
\end{itemize}

\subsection{Unterschiede}

\begin{itemize}
    \item \textbf{Datenquelle:}
    \begin{itemize}
        \item \texttt{extract\_skills\_from\_long\_note.py}: Freitextfeld \texttt{long\_note} in \texttt{candidates}.
        \item \texttt{extract\_skills\_from\_cv\_pdf.py}: PDF in \texttt{cv\_pdf} (BYTEA) in \texttt{candidates}.
        \item \texttt{job\_skills\_extraktion.py}: Stellenbeschreibung \texttt{job\_description} und
              \texttt{long\_note} in \texttt{jobs}.
    \end{itemize}
    \item \textbf{Zielspalte:}
    \begin{itemize}
        \item Kandidaten: \texttt{skills} in der Tabelle \texttt{candidates}.
        \item Jobs: \texttt{sonstiges\_anforderungen} in der Tabelle \texttt{jobs}.
    \end{itemize}
    \item \textbf{Granularität:}
    \begin{itemize}
        \item Die Kandidaten-Skripte arbeiten interaktiv pro Einzel-ID.
        \item Das Job-Skript verarbeitet in einem Lauf viele Datensätze
              und fasst Updates zusammen.
    \end{itemize}
\end{itemize}

\section{Grenzen und mögliche Erweiterungen}

\begin{itemize}
    \item \textbf{Abhängigkeit von der Skill-Liste:} Nur Skills, die in
          \texttt{data/skills.csv} vorkommen, können erkannt werden. Neue
          Skills müssen dort ergänzt werden.
    \item \textbf{Kontextlose Erkennung:} Die Logik prüft aktuell nur auf
          Vorkommen von Begriffen, nicht auf inhaltlichen Kontext
          (z.B. ``kein Python'' wird trotzdem als Python-Skill erkannt).
    \item \textbf{Fuzzy-Schwellwert:} Der Standardwert von 85 ist ein
          Kompromiss zwischen Empfindlichkeit und Präzision. Für andere
          Anwendungsfälle könnte eine Anpassung sinnvoll sein.
    \item \textbf{Mehrsprachigkeit:} Die aktuelle Liste ist auf deutsche
          Bezeichnungen ausgelegt. Für mehrsprachige Lebensläufe könnten
          zusätzliche Spalten (z.B. englische Namen) ergänzt werden.
\end{itemize}

\end{document}
